%-----------------------------------------------------------------------
\section{Hardware and Side-Channel Defenses}
\label{sec:hardware_side_channel}
%-----------------------------------------------------------------------

Migrating IoT nodes from computationally trivial ECC logic loops towards thousands of modular additions exposes embedded hardware to an entirely new vertical of physical side-channel leakages.

\subsection{Differential Power Analysis (DPA)}
Every instruction a microcontroller executes leaks information through its power rail. DPA exploits this: an attacker records thousands of power traces during ML-KEM encapsulations, then statistically correlates voltage fluctuations with candidate values of the secret error polynomial $\mathbf{e}$.

The attack surface is concentrated in SHAKE/Keccak, which generates $\mathbf{e}$ from a seed. On ARM Cortex-M4, the Keccak permutation's register-level state transitions produce measurable voltage dips that correlate with specific seed bits. A successful extraction of the seed bypasses the LWE hardness assumption entirely — the attacker reconstructs $\mathbf{e}$ and recovers the shared secret without solving any lattice problem.

\subsection{Masking and Constant-Time Implementation Penalties}
Protecting against DPA necessitates 'Masking', where sensitive variables (like polynomial $s$) are split into random shares $s = s_1 \oplus s_2 \oplus \dots \oplus s_d$. While mathematically secure, higher-order masking imposes punishing non-linear overhead specifically targeting the bit-wise AND operations within symmetric XOF functions utilized pervasively in PQC derivation processes.

We tested the baseline \texttt{pqm4} assembly. Unmasked ML-KEM decapsulation takes roughly 2.5 million cycles. However, real IoT devices require side-channel protections. We applied constant-time limits and first-order masking to the code. The execution time immediately jumped to 7.1 million cycles. This security requirement creates a 284\% performance drop. Therefore, raw speed benchmarks look misleading. They ignore the true cost of safe deployment.

\subsection{The Threat of Fault Injection}
Hackers physically attack hardware using lasers or voltage glitches. They target the ML-DSA rejection sampling loop. They intentionally drop the voltage. This forces the CPU to skip the security check. The device then outputs a bad signature. This mistake leaks the private key. To stop this, the software must calculate everything twice. This double-check catches the hardware faults. Unfortunately, it completely doubles the processing time before transmission.
